\begin{titlepage}
\newgeometry{margin=0.85in}
\thispagestyle{empty}

\begin{center}
{\large University of Rochester\\}
{\normalsize Warner School of Education\\[0.35in]}

{\color{PrimaryRed}\rule{\textwidth}{1.4pt}}\\[0.22in]
{\Huge\bfseries Designing a Minecraft World\\[0.08in] That Carries Identity and\\[0.08in] Community Knowledge}\\[0.18in]
{\Large Mini-Unit Packet}\\[0.22in]
{\color{PrimaryRed}\rule{0.78\textwidth}{0.7pt}}\\[0.32in]

\colorbox{SoftBlue}{%
\parbox{0.88\textwidth}{%
\centering\vspace{0.12in}
\textbf{Layered Genius: Designing for Joy, Identity, Skills, Intellect, and Criticality}\\[0.05in]
Grade 4--5 Computer Science / Digital Literacy mini-unit centered on Minecraft Education, multimodal meaning-making, identity-safe design, debugging, disciplinary writing, intellectual analysis, and critical digital representation
\vspace{0.12in}
}%
}
\end{center}

\vspace{0.38in}

\begin{center}
\renewcommand{\arraystretch}{1.15}
\setlength{\tabcolsep}{4pt}
\begin{tabular}{>{\bfseries\RaggedLeft}p{1.65in} p{4.9in}}
Student & Piter Zacari Garcia Bautista \\
Course & EDU498: Literacy Learning as Social Practice \\
Session Focus & Toward the Pursuit of Criticality \\
Assignment & CHRE Mini-Unit Extension \\
Mini-Unit Focus & Minecraft world-building as identity, community, coding literacy, disciplinary writing, intellectual design analysis, and critical digital representation \\
Student Profile & Ethan, Student 5, Computer Science \\
Submission Date & June 2026 \\
\end{tabular}
\end{center}

\vspace{0.34in}

\begin{center}
\begin{minipage}{0.86\textwidth}
\itshape
This revision extends the mini-unit through the pursuit of Criticality. It treats Minecraft Education as a disciplinary literacy space where students build, code, debug, document, explain, analyze, and question how digital design choices can represent, misrepresent, include, exclude, preserve, or erase community stories.

\vspace{0.22in}

\scriptsize\centering\textbf{Revision note:} New additions are shown in \textcolor{LessonTwoAccent}{blue}. This version keeps Joy, Identity, Skills, and Intellect developed while deepening Criticality through questions of representation, access, power, respectful design, and repair.
\end{minipage}
\end{center}

\restoregeometry
\end{titlepage}
\clearpage
\onehalfspacing

\section*{Mini-Unit Snapshot}

\textbf{Task:} extend the mini-unit to include the pursuit of Criticality by asking students to analyze how digital design choices communicate meaning and to question how digital spaces can include, exclude, preserve, erase, flatten, or repair community stories.

\textbf{Grade / discipline:} Grade 4--5 Computer Science / Digital Literacy

\textbf{Student profile:} Ethan, a Seneca / Haudenosaunee and mixed-identity student who enjoys coding, game design, robotics, collaboration, outdoor learning, and community storytelling.

\textbf{Mini-unit focus:} Students use Minecraft Education to plan and build a small digital place that communicates identity, community knowledge, belonging, or future possibility while critically examining representation and access in digital spaces.

\textbf{Essential question:} How can code, design, and digital place-making communicate who we are, what matters to us, what our communities know, and what needs to be represented more respectfully?

\textbf{End goal/backward design target:} \textcolor{LessonTwoAccent}{By the end of the unit, students will share a Minecraft place and a short Developer Design Note explaining how specific build/code choices communicate meaning for an audience. They will also identify one representation or access issue, revise one feature, and explain how that revision makes the world clearer, more respectful, more welcoming, or more accessible.}

\textbf{Unit context and purpose:} This three-lesson unit is designed for a Grade 4--5 Computer Science / Digital Literacy setting using Minecraft Education. Students use digital world-building as a literacy practice: they plan, code, debug, label, explain, analyze, question, and revise a digital place that carries meaning. The purpose is to help students see coding as authorship, identity, community knowledge, audience-aware design, and critical meaning-making. The unit foregrounds Joy and Identity by giving students choice, visible progress, and a protected way to connect design with who they are or what they value. It strengthens Skills through planning, sequencing, debugging, documentation, and presentation. It develops Intellect by asking students to analyze how code, signs, paths, colors, structures, repeated patterns, accessibility supports, and audience choices shape how others read a digital world. \textcolor{LessonTwoAccent}{It deepens Criticality by asking students to question whose stories appear in digital worlds, whose knowledge is missing, who benefits from certain representations, who may be harmed, and how designers can use technology to repair rather than reproduce exclusion.}

\textbf{Three-lesson Arc:}
\begin{itemize}[leftmargin=*]
    \item \textcolor{LessonOneAccent}{\textbf{Lesson 1: Plan a Meaningful Place.}} Students compare a generic build with a meaningful digital place, choose a place concept, sketch meaningful features, decompose the build into steps, and begin asking whose stories are represented or missing in digital spaces.
    \item \textcolor{LessonTwoAccent}{\textbf{Lesson 2: Build, Code, Debug, and Analyze.}} Students use Minecraft / MakeCode to build, test, revise, and document one coded or repeated feature that helps the world communicate meaning while considering whether the design is respectful and accessible.
    \item \textcolor{LessonThreeAccent}{\textbf{Lesson 3: Share, Revise, and Reflect on Representation.}} Students present the world, submit a Developer Design Note, explain design choices, receive feedback, revise one feature for clarity, welcome, or accessibility, and reflect on representation, missing stories, and repair.
\end{itemize}

\textcolor{LessonTwoAccent}{Across the arc, the disciplinary thread stays consistent: students move from idea to plan, from plan to build/code, from code to debugging evidence, and from debugging evidence to a critical explanation of how design communicates meaning and affects others. This makes computer science skill visible without separating it from identity, joy, access, community knowledge, and justice.}

\newpage
\section*{Source Trail and Rationale for the Pick}

This work grows from my work toward ``\emph{Layered Genius: Designing for Joy and Identity}'' and from Muhammad's Identity pursuit, which frames identity as self-definition. It also carries forward my joy commitment of ``belonging before building'' and my placement observations of debugging, collaborative work, peer teaching, and student-led problem framing. I also draw on Moje, Luke, Davies, and Street's literacy-and-identity lens, Skerrett's attention to repositioning students as capable contributors, and Moje's disciplinary literacy framework. These sources help name coding, debugging, design, and explanation as literacy practices. Minecraft Education/MakeCode and the CSTA/ISTE standards serve as practical anchors for sequencing, debugging, computational thinking, creative communication, collaboration, and multimodal digital design.

\textcolor{LessonTwoAccent}{For this revision, I also use the Class 10 pursuit of Criticality to move the mini-unit beyond representation as a surface feature. Students are asked to read digital worlds as designed texts: who is centered, who is missing, what messages are repeated, what forms of access are assumed, and how technology can be redesigned toward repair. This connects to Freire's idea of reading the world and to Muhammad's pursuit of Criticality because students are not only noticing injustice; they are studying how it is produced and imagining how design can interrupt it.}

\section*{Student Profile and Identity Safety}

Ethan's profile matters because the lesson can recognize him as a coder, designer, storyteller, and community-minded problem solver. At the same time, the lesson should honor his identity without requiring him to explain, display, or represent his culture for the class. Instead, it offers identity-safe choice by allowing students to build from lived experience, imagination, community, interests, family, land, language, hobbies, or future hopes. Because Ethan's profile includes Seneca / Haudenosaunee heritage, culturally specific symbols or stories should be student-chosen and handled with care. Private disclosure is never required. Fictionalized, symbolic, or aspirational places are valid.

\textbf{Identity design move:} The lesson treats identity as a design resource, not a disclosure requirement. Students choose what counts as meaningful, and the teacher assesses how clearly design choices communicate meaning. \textcolor{LessonTwoAccent}{The criticality layer protects student agency by asking students to examine representation and access without forcing personal or cultural disclosure.}

\section*{Discipline Standards, Revised Pursuit, and Learning Activity Ideas}

\textcolor{LessonTwoAccent}{\textbf{Learning standard CSTA Computer Science:}}

\textcolor{LessonTwoAccent}{\textbf{1B-AP-17:} Describe choices made during program development using code comments, presentations, and demonstrations.}

\textcolor{LessonTwoAccent}{\textbf{Revised:} Students will design, evaluate, and refine a Minecraft Education world that communicates identity, community knowledge, belonging, or future possibility. Students will analyze how their code, signs, paths, structures, repeated patterns, and audience supports shape how others understand the digital space. They will also question how digital spaces can represent, misrepresent, include, or exclude people and communities, then revise one feature toward access, respect, or repair.}

\textcolor{LessonTwoAccent}{\textbf{Digital Representation and Access to Education:} Addressing how digital spaces can include or exclude students' identities, community knowledge, and access needs while promoting respectful and accessible design. Example: Creating a Minecraft Education world that represents a meaningful place, community story, or future possibility without requiring students to disclose private cultural or personal information. Students use critical questions to examine whose stories are visible, whose stories are missing, and how design choices can repair or reproduce exclusion.}

\textcolor{LessonTwoAccent}{\textbf{Disciplinary literacy move:} In this mini-unit, writing is not separate from computer science. Students write like beginning designers/developers: they state a design goal, name steps, describe a bug or revision, explain how a feature works, connect technical choices to audience meaning, and reflect on representation, access, and repair.}

\section*{Pursuits Alignment}

\renewcommand{\arraystretch}{1.08}
\setlength{\tabcolsep}{3pt}
\begin{tabularx}{\textwidth}{|>{\raggedright\arraybackslash}p{0.09\textwidth}|Y|}
\hline
\rowcolor{SoftBlue}
\textbf{Pursuit} & \textbf{How to Address it} \\ \hline
\rowcolor{JoySoft}
Joy & Students experience choice, visible progress, peer support, debugging as discovery, and the satisfaction of building something shareable. \\ \hline
\rowcolor{IdentitySoft}
Identity & Students design a meaningful place and explain how specific features communicate belonging, community, memory, culture, interests, land, imagination, or future possibility with dignity. \\ \hline
\rowcolor{SkillsSoft}
Skills & Students practice sequencing, decomposition, repeated patterns, debugging, reasoning, screenshot use, partner communication, and oral/written explanation. They also complete a Developer Design Note that documents their design goal, steps, bug/revision, and audience meaning. \\ \hline
\rowcolor{IntellectSoft}
Intellect & Students analyze how code, space, signs, paths, color, structures, repeated patterns, interaction, accessibility supports, and audience choices work together to create meaning in a digital world. They evaluate whether their world communicates clearly, revise one feature based on feedback, and explain how the design helps an audience understand identity, community knowledge, belonging, or future possibility. \\ \hline
\rowcolor{CriticalitySoft}
Criticality & \textcolor{LessonTwoAccent}{Students question whose stories are represented in digital worlds, whose knowledge is missing, how technology can preserve or erase community stories, and how designers can avoid using culture, identity, or community knowledge as decoration. Students also ask who benefits from certain representations, who may be harmed, and how technology can be redesigned toward access, respect, repair, and transformation.} \\ \hline
\end{tabularx}

\section*{Multimodal and Multiliteracies Design}

Minecraft functions here as a text students author. They read digital space, sketch a plan, build visually, use procedural thinking, explain orally, label in writing, analyze design choices, question representation, and revise through peer feedback.

\textbf{Modes included:} spatial design, visual design, coding/sequencing, oral explanation, written labels or reflection, screenshots, collaboration, peer feedback, audience-aware revision, and critical questioning.

\textbf{Access moves:} sketch-first or build-first choice; projected model; partner roles; sentence stems; optional oral explanation; screenshots as memory support; chunked directions; low-floor/high-ceiling challenge options; and a small revision target instead of requiring students to rebuild the whole world.

\textbf{Why this supports multiliteracies:} Students are not only writing or answering questions; they are composing meaning across space, image, code, talk, labels, screenshots, and peer feedback. A path, sign, repeated pattern, gathering space, accessibility feature, or coded sequence can all function as evidence of meaning-making.

\textcolor{LessonTwoAccent}{\textbf{Critical literacy layer:} Students read their world as designers and as critical users of technology. They ask what an audience can infer from paths, signs, repeated patterns, coded movement, points of interaction, and accessibility supports. They also ask whose stories are centered, whose stories are missing, and how one design revision can make the world more respectful, accessible, or reparative.}
